% Options for packages loaded elsewhere
\PassOptionsToPackage{unicode}{hyperref}
\PassOptionsToPackage{hyphens}{url}
\documentclass[
]{article}
\usepackage{xcolor}
\usepackage[margin=1in]{geometry}
\usepackage{amsmath,amssymb}
\setcounter{secnumdepth}{-\maxdimen} % remove section numbering
\usepackage{iftex}
\ifPDFTeX
  \usepackage[T1]{fontenc}
  \usepackage[utf8]{inputenc}
  \usepackage{textcomp} % provide euro and other symbols
\else % if luatex or xetex
  \usepackage{unicode-math} % this also loads fontspec
  \defaultfontfeatures{Scale=MatchLowercase}
  \defaultfontfeatures[\rmfamily]{Ligatures=TeX,Scale=1}
\fi
\usepackage{lmodern}
\ifPDFTeX\else
  % xetex/luatex font selection
\fi
% Use upquote if available, for straight quotes in verbatim environments
\IfFileExists{upquote.sty}{\usepackage{upquote}}{}
\IfFileExists{microtype.sty}{% use microtype if available
  \usepackage[]{microtype}
  \UseMicrotypeSet[protrusion]{basicmath} % disable protrusion for tt fonts
}{}
\makeatletter
\@ifundefined{KOMAClassName}{% if non-KOMA class
  \IfFileExists{parskip.sty}{%
    \usepackage{parskip}
  }{% else
    \setlength{\parindent}{0pt}
    \setlength{\parskip}{6pt plus 2pt minus 1pt}}
}{% if KOMA class
  \KOMAoptions{parskip=half}}
\makeatother
\usepackage{color}
\usepackage{fancyvrb}
\newcommand{\VerbBar}{|}
\newcommand{\VERB}{\Verb[commandchars=\\\{\}]}
\DefineVerbatimEnvironment{Highlighting}{Verbatim}{commandchars=\\\{\}}
% Add ',fontsize=\small' for more characters per line
\usepackage{framed}
\definecolor{shadecolor}{RGB}{248,248,248}
\newenvironment{Shaded}{\begin{snugshade}}{\end{snugshade}}
\newcommand{\AlertTok}[1]{\textcolor[rgb]{0.94,0.16,0.16}{#1}}
\newcommand{\AnnotationTok}[1]{\textcolor[rgb]{0.56,0.35,0.01}{\textbf{\textit{#1}}}}
\newcommand{\AttributeTok}[1]{\textcolor[rgb]{0.13,0.29,0.53}{#1}}
\newcommand{\BaseNTok}[1]{\textcolor[rgb]{0.00,0.00,0.81}{#1}}
\newcommand{\BuiltInTok}[1]{#1}
\newcommand{\CharTok}[1]{\textcolor[rgb]{0.31,0.60,0.02}{#1}}
\newcommand{\CommentTok}[1]{\textcolor[rgb]{0.56,0.35,0.01}{\textit{#1}}}
\newcommand{\CommentVarTok}[1]{\textcolor[rgb]{0.56,0.35,0.01}{\textbf{\textit{#1}}}}
\newcommand{\ConstantTok}[1]{\textcolor[rgb]{0.56,0.35,0.01}{#1}}
\newcommand{\ControlFlowTok}[1]{\textcolor[rgb]{0.13,0.29,0.53}{\textbf{#1}}}
\newcommand{\DataTypeTok}[1]{\textcolor[rgb]{0.13,0.29,0.53}{#1}}
\newcommand{\DecValTok}[1]{\textcolor[rgb]{0.00,0.00,0.81}{#1}}
\newcommand{\DocumentationTok}[1]{\textcolor[rgb]{0.56,0.35,0.01}{\textbf{\textit{#1}}}}
\newcommand{\ErrorTok}[1]{\textcolor[rgb]{0.64,0.00,0.00}{\textbf{#1}}}
\newcommand{\ExtensionTok}[1]{#1}
\newcommand{\FloatTok}[1]{\textcolor[rgb]{0.00,0.00,0.81}{#1}}
\newcommand{\FunctionTok}[1]{\textcolor[rgb]{0.13,0.29,0.53}{\textbf{#1}}}
\newcommand{\ImportTok}[1]{#1}
\newcommand{\InformationTok}[1]{\textcolor[rgb]{0.56,0.35,0.01}{\textbf{\textit{#1}}}}
\newcommand{\KeywordTok}[1]{\textcolor[rgb]{0.13,0.29,0.53}{\textbf{#1}}}
\newcommand{\NormalTok}[1]{#1}
\newcommand{\OperatorTok}[1]{\textcolor[rgb]{0.81,0.36,0.00}{\textbf{#1}}}
\newcommand{\OtherTok}[1]{\textcolor[rgb]{0.56,0.35,0.01}{#1}}
\newcommand{\PreprocessorTok}[1]{\textcolor[rgb]{0.56,0.35,0.01}{\textit{#1}}}
\newcommand{\RegionMarkerTok}[1]{#1}
\newcommand{\SpecialCharTok}[1]{\textcolor[rgb]{0.81,0.36,0.00}{\textbf{#1}}}
\newcommand{\SpecialStringTok}[1]{\textcolor[rgb]{0.31,0.60,0.02}{#1}}
\newcommand{\StringTok}[1]{\textcolor[rgb]{0.31,0.60,0.02}{#1}}
\newcommand{\VariableTok}[1]{\textcolor[rgb]{0.00,0.00,0.00}{#1}}
\newcommand{\VerbatimStringTok}[1]{\textcolor[rgb]{0.31,0.60,0.02}{#1}}
\newcommand{\WarningTok}[1]{\textcolor[rgb]{0.56,0.35,0.01}{\textbf{\textit{#1}}}}
\usepackage{graphicx}
\makeatletter
\newsavebox\pandoc@box
\newcommand*\pandocbounded[1]{% scales image to fit in text height/width
  \sbox\pandoc@box{#1}%
  \Gscale@div\@tempa{\textheight}{\dimexpr\ht\pandoc@box+\dp\pandoc@box\relax}%
  \Gscale@div\@tempb{\linewidth}{\wd\pandoc@box}%
  \ifdim\@tempb\p@<\@tempa\p@\let\@tempa\@tempb\fi% select the smaller of both
  \ifdim\@tempa\p@<\p@\scalebox{\@tempa}{\usebox\pandoc@box}%
  \else\usebox{\pandoc@box}%
  \fi%
}
% Set default figure placement to htbp
\def\fps@figure{htbp}
\makeatother
\setlength{\emergencystretch}{3em} % prevent overfull lines
\providecommand{\tightlist}{%
  \setlength{\itemsep}{0pt}\setlength{\parskip}{0pt}}
\usepackage{bookmark}
\IfFileExists{xurl.sty}{\usepackage{xurl}}{} % add URL line breaks if available
\urlstyle{same}
\hypersetup{
  pdftitle={White-Box vs.~Black-Box: Weibull Estimation in R},
  pdfauthor={Reliability Analysis Project},
  hidelinks,
  pdfcreator={LaTeX via pandoc}}

\title{White-Box vs.~Black-Box: Weibull Estimation in R}
\author{Reliability Analysis Project}
\date{2026-04-27}

\begin{document}
\maketitle

{
\setcounter{tocdepth}{2}
\tableofcontents
}
\section{Introduction}\label{introduction}

In the main \texttt{02\_estimation\_methods.R} script, we implemented
\textbf{``White-Box''} code. This means we manually defined the
underlying mathematics (like explicitly writing the log-likelihood
functions, algebra for Bergman Weights, and root-finding for Moments).
We did this to create an authentic educational pipeline where students
can see exactly \emph{how} the algorithms arrive at their parameters.

However, in a production environment where you just want results
quickly, R provides \textbf{``Black-Box''} functions. These are
optimized, robust, one-liner functions embedded in standard packages
that solve the same problems instantly.

This document contrasts the verbosity of our educational code against
the brevity of production-ready packages.

\begin{center}\rule{0.5\linewidth}{0.5pt}\end{center}

\section{1. Maximum Likelihood Estimation
(MLE)}\label{maximum-likelihood-estimation-mle}

MLE is the industry standard for estimating Weibull parameters with
right-censored data.

\subsubsection{The Detailed ``White-Box''
Approach}\label{the-detailed-white-box-approach}

In our code, we explicitly mapped the failure likelihoods against the
censored survival probabilities, resulting in an \textasciitilde30-line
custom implementation passed to a general-purpose optimizer
(\texttt{optim}):

\begin{Shaded}
\begin{Highlighting}[]
\CommentTok{\# Manual Log{-}Likelihood \& Nelder{-}Mead Optimization (\textasciitilde{}30 Lines)}
\NormalTok{neg\_loglik }\OtherTok{\textless{}{-}} \ControlFlowTok{function}\NormalTok{(log\_params) \{}
\NormalTok{  k   }\OtherTok{\textless{}{-}} \FunctionTok{exp}\NormalTok{(log\_params[}\DecValTok{1}\NormalTok{])}
\NormalTok{  lam }\OtherTok{\textless{}{-}} \FunctionTok{exp}\NormalTok{(log\_params[}\DecValTok{2}\NormalTok{])}
  
\NormalTok{  t\_ratio }\OtherTok{\textless{}{-}}\NormalTok{ t }\SpecialCharTok{/}\NormalTok{ lam}
  
  \CommentTok{\# Log{-}PDF for failures}
\NormalTok{  ll\_fail }\OtherTok{\textless{}{-}} \FunctionTok{sum}\NormalTok{(event }\SpecialCharTok{*}\NormalTok{ (}\FunctionTok{log}\NormalTok{(k) }\SpecialCharTok{{-}} \FunctionTok{log}\NormalTok{(lam) }\SpecialCharTok{+}\NormalTok{ (k }\SpecialCharTok{{-}} \DecValTok{1}\NormalTok{) }\SpecialCharTok{*} \FunctionTok{log}\NormalTok{(t\_ratio) }\SpecialCharTok{{-}}\NormalTok{ t\_ratio}\SpecialCharTok{\^{}}\NormalTok{k))}
  \CommentTok{\# Log{-}survival for censored}
\NormalTok{  ll\_cens }\OtherTok{\textless{}{-}} \FunctionTok{sum}\NormalTok{((}\DecValTok{1} \SpecialCharTok{{-}}\NormalTok{ event) }\SpecialCharTok{*}\NormalTok{ (}\SpecialCharTok{{-}}\NormalTok{t\_ratio}\SpecialCharTok{\^{}}\NormalTok{k))}
  
  \FunctionTok{return}\NormalTok{(}\SpecialCharTok{{-}}\NormalTok{(ll\_fail }\SpecialCharTok{+}\NormalTok{ ll\_cens))}
\NormalTok{\}}

\NormalTok{result }\OtherTok{\textless{}{-}} \FunctionTok{optim}\NormalTok{(}\AttributeTok{par =} \FunctionTok{c}\NormalTok{(}\FunctionTok{log}\NormalTok{(}\FloatTok{1.5}\NormalTok{), }\FunctionTok{log}\NormalTok{(mean\_t)), }
                \AttributeTok{fn =}\NormalTok{ neg\_loglik, }
                \AttributeTok{method =} \StringTok{"Nelder{-}Mead"}\NormalTok{)}

\NormalTok{k\_hat   }\OtherTok{\textless{}{-}} \FunctionTok{exp}\NormalTok{(result}\SpecialCharTok{$}\NormalTok{par[}\DecValTok{1}\NormalTok{])}
\NormalTok{lam\_hat }\OtherTok{\textless{}{-}} \FunctionTok{exp}\NormalTok{(result}\SpecialCharTok{$}\NormalTok{par[}\DecValTok{2}\NormalTok{])}
\end{Highlighting}
\end{Shaded}

\subsubsection{The ``Black-Box''
One-Liner}\label{the-black-box-one-liner}

In production, you would never write your own likelihood function. You
would use the \texttt{survival} package, which is arguably the most
battle-tested survival analysis library in existence.

\begin{Shaded}
\begin{Highlighting}[]
\CommentTok{\# Using the survival package (1 Line)}
\FunctionTok{library}\NormalTok{(survival)}

\NormalTok{fit }\OtherTok{\textless{}{-}} \FunctionTok{survreg}\NormalTok{(}\FunctionTok{Surv}\NormalTok{(lifetime\_days, failed) }\SpecialCharTok{\textasciitilde{}} \DecValTok{1}\NormalTok{, }\AttributeTok{data =}\NormalTok{ drive\_data, }\AttributeTok{dist =} \StringTok{"weibull"}\NormalTok{)}

\CommentTok{\# Note: survreg uses the "Extreme Value" parameterization. }
\CommentTok{\# To convert to standard Weibull shape/scale:}
\NormalTok{k\_hat   }\OtherTok{\textless{}{-}} \DecValTok{1} \SpecialCharTok{/}\NormalTok{ fit}\SpecialCharTok{$}\NormalTok{scale}
\NormalTok{lam\_hat }\OtherTok{\textless{}{-}} \FunctionTok{exp}\NormalTok{(}\FunctionTok{coef}\NormalTok{(fit))}
\end{Highlighting}
\end{Shaded}

\emph{Code Reduction:} \textbf{30 lines} down to \textbf{1 line}.

\begin{center}\rule{0.5\linewidth}{0.5pt}\end{center}

\section{2. OLS / Median Rank Regression
(MRR)}\label{ols-median-rank-regression-mrr}

OLS fits a linear regression line to coordinate transformations of the
empirical failures.

\subsubsection{The Detailed ``White-Box''
Approach}\label{the-detailed-white-box-approach-1}

We manually calculated Benard's median ranks, transformed the axes to
the Weibull scale, and explicitly applied R's \texttt{lm()} function:

\begin{Shaded}
\begin{Highlighting}[]
\CommentTok{\# Manual Rank Transformation (\textasciitilde{}15 lines)}
\NormalTok{t\_fail }\OtherTok{\textless{}{-}} \FunctionTok{sort}\NormalTok{(drive\_data}\SpecialCharTok{$}\NormalTok{lifetime\_days[drive\_data}\SpecialCharTok{$}\NormalTok{failed }\SpecialCharTok{==} \DecValTok{1}\NormalTok{])}
\NormalTok{n }\OtherTok{\textless{}{-}} \FunctionTok{length}\NormalTok{(t\_fail)}

\CommentTok{\# Benard\textquotesingle{}s formula}
\NormalTok{i }\OtherTok{\textless{}{-}} \FunctionTok{seq\_len}\NormalTok{(n)}
\NormalTok{F\_i }\OtherTok{\textless{}{-}}\NormalTok{ (i }\SpecialCharTok{{-}} \FloatTok{0.3}\NormalTok{) }\SpecialCharTok{/}\NormalTok{ (n }\SpecialCharTok{+} \FloatTok{0.4}\NormalTok{)}

\CommentTok{\# Linearization}
\NormalTok{Y }\OtherTok{\textless{}{-}} \FunctionTok{log}\NormalTok{(}\SpecialCharTok{{-}}\FunctionTok{log}\NormalTok{(}\DecValTok{1} \SpecialCharTok{{-}}\NormalTok{ F\_i)) }
\NormalTok{X }\OtherTok{\textless{}{-}} \FunctionTok{log}\NormalTok{(t\_fail)}

\NormalTok{fit }\OtherTok{\textless{}{-}} \FunctionTok{lm}\NormalTok{(Y }\SpecialCharTok{\textasciitilde{}}\NormalTok{ X)}
\NormalTok{k\_hat   }\OtherTok{\textless{}{-}} \FunctionTok{unname}\NormalTok{(}\FunctionTok{coef}\NormalTok{(fit)[}\DecValTok{2}\NormalTok{])}
\NormalTok{lam\_hat }\OtherTok{\textless{}{-}} \FunctionTok{exp}\NormalTok{(}\SpecialCharTok{{-}}\FunctionTok{unname}\NormalTok{(}\FunctionTok{coef}\NormalTok{(fit)[}\DecValTok{1}\NormalTok{]) }\SpecialCharTok{/}\NormalTok{ k\_hat)}
\end{Highlighting}
\end{Shaded}

\subsubsection{The ``Black-Box''
One-Liner}\label{the-black-box-one-liner-1}

If you need Rank Regression in R, the specialized \texttt{WeibullR}
package handles all of this automatically, converting and plotting in
one swift move.

\begin{Shaded}
\begin{Highlighting}[]
\CommentTok{\# Using WeibullR (1 Line)}
\FunctionTok{library}\NormalTok{(WeibullR)}

\NormalTok{fit }\OtherTok{\textless{}{-}} \FunctionTok{wblr.fit}\NormalTok{(}\FunctionTok{wblr}\NormalTok{(t\_fail), }\AttributeTok{method.fit =} \StringTok{"rr"}\NormalTok{)}
\end{Highlighting}
\end{Shaded}

\emph{Code Reduction:} \textbf{15 lines} down to \textbf{1 line}.

\begin{center}\rule{0.5\linewidth}{0.5pt}\end{center}

\section{3. Kaplan-Meier Survival
Benchmark}\label{kaplan-meier-survival-benchmark}

The Kaplan-Meier estimator is exactly what we just added to the lecture
documents to prove why MLE works better than WLS for censored drives.

\subsubsection{The Detailed ``White-Box''
Approach}\label{the-detailed-white-box-approach-2}

Plotting this manually requires grouping data, keeping track of the risk
set at each time step \(t\), calculating running survival probabilities
\(\prod (1 - d_i/n_i)\), and building a step-function graph. If coded
sequentially, it easily takes \textbf{20-40 lines}.

\subsubsection{The ``Black-Box''
One-Liner}\label{the-black-box-one-liner-2}

The \texttt{survival} package calculates this using the formula
instantly.

\begin{Shaded}
\begin{Highlighting}[]
\CommentTok{\# Using survival and broom for instant calculation (\textasciitilde{}2 Lines)}
\FunctionTok{library}\NormalTok{(survival)}

\NormalTok{km\_fit }\OtherTok{\textless{}{-}} \FunctionTok{survfit}\NormalTok{(}\FunctionTok{Surv}\NormalTok{(lifetime\_days, failed) }\SpecialCharTok{\textasciitilde{}} \DecValTok{1}\NormalTok{, }\AttributeTok{data =}\NormalTok{ drive\_data)}
\CommentTok{\# Plotting is as easy as:}
\FunctionTok{plot}\NormalTok{(km\_fit, }\AttributeTok{conf.int =} \ConstantTok{TRUE}\NormalTok{, }\AttributeTok{mark.time =} \ConstantTok{TRUE}\NormalTok{)}
\end{Highlighting}
\end{Shaded}

\begin{center}\rule{0.5\linewidth}{0.5pt}\end{center}

\section{Conclusion}\label{conclusion}

While the \textbf{White-Box} code makes our project significantly
longer, it acts as a digital textbook. It strips away the magic of
libraries like \texttt{survival} and \texttt{WeibullR} to reveal exactly
how statistical mathematics translates to code.

When you move into actual corporate analytics, you will discard the
detailed scripts found in \texttt{02\_estimation\_methods.R} and instead
pipeline everything through the \textbf{Black-Box} one-liners shown
above!

\end{document}
